% ============================================================================
%  CAPÍTULO III: DESARROLLO DEL TRABAJO
% ============================================================================

\section{Desarrollo del trabajo}

Este capítulo documenta el proceso completo de instalación, configuración y
prueba del servidor DHCP implementado con Kea sobre el sistema Debian del
proyecto. Se describe desde la instalación de los paquetes necesarios hasta
las pruebas de asignación automática de direcciones IP realizadas desde
equipos clientes.

\section{Servidor DHCP}

El servidor DHCP se implementó utilizando Kea, el servidor DHCP de ISC
(\textit{Internet Systems Consortium}) y sucesor del clásico
\texttt{dhcpd}, seleccionado por su soporte activo, su configuración basada
en JSON y su manejo eficiente de las concesiones \cite{keaDocs}. A
continuación se describe el proceso de implementación.

% -------------------------------------------------------
\subsection{Arquitectura del sistema}

El servicio se desplegó de manera centralizada sobre el servidor Debian de la
red local \texttt{192.168.0.0/24}. Para que el servidor DHCP sea el único que
otorgue direcciones, se desactivó el servicio DHCP del router de la red,
evitando así que dos servidores respondan a las peticiones \textsc{Discover}
de los clientes con ofertas distintas.

La red se dividió en dos segmentos lógicos según la política del proyecto:

\begin{itemize}
  \item \textbf{Equipos de rol (192.168.0.1--192.168.0.109):} incluye al
        router, al propio servidor Debian (\texttt{192.168.0.105}) y a todos
        los servicios (DNS, web, archivos, correo) que requieren una dirección
        IP fija.
  \item \textbf{Equipos de usuario (192.168.0.110--192.168.0.254):} rango de
        direcciones que el servidor DHCP puede entregar dinámicamente a los
        dispositivos de los usuarios de la red.
\end{itemize}

El servidor DHCP entrega a cada cliente, junto con la dirección IP, la
máscara de subred (\texttt{255.255.255.0}), la puerta de enlace
(\texttt{192.168.0.1}), el servidor DNS del proyecto
(\texttt{192.168.0.105}) y el dominio de búsqueda (\texttt{sudoers.lan}).
De esta manera, el cliente queda configurado sin intervención manual y puede
resolver los nombres de los servicios internos inmediatamente.

La figura \ref{fig:arquitectura-dhcp} presenta la arquitectura del servicio
implementado.

\figuraAPA{fig:arquitectura-dhcp}{Arquitectura del servicio DHCP: clientes, router y servidor}{%
  \begin{tikzpicture}[
      node distance=1.2cm and 1.6cm,
      ser/.style={draw, rounded corners, minimum width=2.6cm,
                  minimum height=0.9cm, align=center, font=\small},
      cli/.style={draw, minimum width=2.4cm, minimum height=0.8cm,
                  align=center, font=\small},
      flujo/.style={-Stealth, thick}
    ]
    \node[ser] (kea) {Kea DHCP\\192.168.0.105};
    \node[cli, left=of kea, yshift=0.7cm] (c1) {Usuario 1\\192.168.0.110+};
    \node[cli, left=of kea, yshift=-0.7cm] (c2) {Usuario 2\\192.168.0.110+};
    \node[ser, right=of kea] (router) {Router\\192.168.0.1\\(DHCP off)};
    \node[ser, right=of router] (dns) {DNS\\192.168.0.105};

    \draw[flujo] (c1.east) -- (kea.west);
    \draw[flujo] (c2.east) -- (kea.west);
    \draw[flujo] (kea.east) -- (router.west);
    \draw[flujo, dashed] (dns.west) -- (router.east);
  \end{tikzpicture}%
}{Elaboración propia.}

% -------------------------------------------------------
\subsection{Instalación del servidor DHCP}

El servidor DHCP se instaló sobre el sistema Debian mediante el administrador
de paquetes APT. Kea se distribuye en paquetes separados por protocolo; para
este proyecto se requiere únicamente el servicio DHCPv4:

\begin{lstlisting}[style=terminal]
$ apt update
$ apt install -y kea-dhcp4-server
\end{lstlisting}

El paquete \texttt{kea-dhcp4-server} instala el binario
\texttt{kea-dhcp4}, crea el usuario de sistema \texttt{\_kea} con el que el
servicio ejecuta y registra la unidad de \texttt{systemd}
\texttt{kea-dhcp4-server.service}. Después de la instalación, la
configuración principal queda en \texttt{/etc/kea/kea-dhcp4.conf}.

% -------------------------------------------------------
\subsection{Configuración del servidor}

La configuración de Kea se define en el archivo
\texttt{/etc/kea/kea-dhcp4.conf}, con formato JSON. A continuación se
describen los bloques más relevantes configurados para el proyecto.

\subsubsection{Interfaz y base de datos de concesiones}

Kea escucha en todas las interfaces del servidor para poder atender los
paquetes \textsc{Discover} que llegan por broadcast desde la LAN. Las
concesiones se almacenan en el archivo
\texttt{/var/lib/kea/kea-leases4.csv}:

\begin{lstlisting}[style=terminal]
"interfaces-config": { "interfaces": ["*"] },
"lease-database": { "type": "memfile", "lfc-interval": 3600 }
\end{lstlisting}

La opción \texttt{lfc-interval} indica cada cuántos segundos se compacta el
archivo de concesiones para mantenerlo eficiente.

\subsubsection{Subred, ámbito y opciones}

Se definió la subred \texttt{192.168.0.0/24} con un único ámbito de
asignación dinámica comprendido entre \texttt{192.168.0.110} y
\texttt{192.168.0.254}:

\begin{lstlisting}[style=terminal]
"subnet4": [{
  "subnet": "192.168.0.0/24",
  "pools": [{ "pool": "192.168.0.110 - 192.168.0.254" }],
  "option-data": [
    { "name": "domain-name-servers", "data": "192.168.0.105" },
    { "name": "domain-name", "data": "sudoers.lan" },
    { "name": "routers", "data": "192.168.0.1" }
  ]
}]
\end{lstlisting}

La separación del rango es intencional: los equipos de rol usan direcciones
inferiores a \texttt{192.168.0.110} (fijas o configuradas fuera del ámbito),
mientras que el ámbito dinámico atiende exclusivamente a los equipos de
usuario. Además, cada cliente recibe automáticamente el servidor DNS del
proyecto (\texttt{192.168.0.105}), el dominio \texttt{sudoers.lan} y la
puerta de enlace \texttt{192.168.0.1}.

% -------------------------------------------------------
\subsection{Desactivación del DHCP del router}

El router de la red contaba con su propio servicio DHCP activo, lo que
provocaría un conflicto: tanto el router como Kea responderían a los
paquetes \textsc{Discover} y los clientes podrían recibir configuraciones
distintas e inconsistentes. Para garantizar que el único servidor DHCP de la
red sea el implementado en el proyecto, se accedió a la interfaz de
administración del router y se deshabilitó el servicio DHCP en la LAN
(DHCP Server - Off), manteniendo únicamente la configuración de la puerta de
enlace y las reglas de red existentes.

Al realizar esta desactivación, los clientes conectados a la red que tuvieran
una IP asignada previamente por el router dependen de renovar su concesión;
una vez que lo hacen, el servidor Kea les asigna una dirección dentro del
ámbito dinámico junto con la configuración de red completa.

% -------------------------------------------------------
\subsection{Verificación del servicio}

Antes de probar con clientes reales se validó la configuración y el estado
del servicio. La sintaxis del archivo de configuración se comprobó con la
opción de prueba de Kea:

\begin{lstlisting}[style=terminal]
$ kea-dhcp4 -t /etc/kea/kea-dhcp4.conf
$ systemctl restart kea-dhcp4-server
$ systemctl status kea-dhcp4-server
\end{lstlisting}

El comando \texttt{kea-dhcp4 -t} devuelve \texttt{Configuration check
successful} si el archivo es válido. Posteriormente se reinició el servicio y
se confirmó su estado activo (\texttt{active (running)}) mediante
\texttt{systemctl}. La figura \ref{fig:captura-estado} muestra el estado del
servicio en ejecución.

\capturaAPA{fig:captura-estado}{Salida de \texttt{systemctl status kea-dhcp4-server} en el servidor}{captura-systemctl-kea.png}{Captura de pantalla de la terminal mostrando el estado activo del servicio \texttt{kea-dhcp4-server} con la ruta de la configuración cargada.}

% -------------------------------------------------------
\subsection{Pruebas de asignación desde clientes}

Para comprobar la asignación automática se utilizó un equipo cliente conectado
a la misma red física. Se ejecutó el cliente DHCP del sistema con
\texttt{dhclient} y se verificó la configuración recibida con \texttt{ip a}:

\begin{lstlisting}[style=terminal]
$ sudo dhclient -v eth0
$ ip a
$ ip route
$ resolvectl status
\end{lstlisting}

El cliente recibió una dirección IP dentro del ámbito dinámico
(\texttt{192.168.0.110}--\texttt{192.168.0.254}), junto con la máscara, la
puerta de enlace y el servidor DNS entregados por Kea. La figura
\ref{fig:captura-cliente} evidencia la asignación en el equipo cliente.

\capturaAPA{fig:captura-cliente}{Dirección IP asignada al cliente con \texttt{ip a} tras el arranque de \texttt{dhclient}}{captura-cliente-ip.png}{Captura de pantalla de la salida de \texttt{ip a} en el cliente mostrando la dirección IP, la máscara \texttt{/24} y el estado de la interfaz (UP).}

% -------------------------------------------------------
\subsection{Registro de concesiones}

Cada asignación queda registrada en el archivo de concesiones de Kea. Su
contenido se puede consultar directamente desde el servidor:

\begin{lstlisting}[style=terminal]
$ cat /var/lib/kea/kea-leases4.csv
\end{lstlisting}

La figura \ref{fig:captura-leases} muestra las concesiones otorgadas por el
servidor, identificando la dirección IP entregada, la dirección MAC del
cliente y el tiempo de validez de cada lease, lo que confirma que la
asignación dinámica funciona correctamente.

\capturaAPA{fig:captura-leases}{Concesiones registradas por Kea en \texttt{/var/lib/kea/kea-leases4.csv}}{captura-leases.png}{Captura de pantalla del archivo \texttt{kea-leases4.csv} mostrando las concesiones activas: dirección IP, hardware address del cliente, hostname y tiempos de renovación/expiracion.}

Finalmente, la conectividad del cliente se validó con \texttt{ping} hacia la
puerta de enlace y hacia el servidor DNS, comprobando que la configuración
entregada por DHCP permite al equipo alcanzar los servicios internos de la
red:

\begin{lstlisting}[style=terminal]
$ ping -c 2 192.168.0.1
$ ping -c 2 192.168.0.105
\end{lstlisting}